\begin{abstract}
针对古代玻璃文物成分数据，构建了关联分析、成分预测与分类鉴定模型。问题一采用卡方检验与Cramér's V系数，发现风化状态与纹饰存在极显著的强关联（$p = 4.7 \times 10^{-5}$, V = 0.5454），而与玻璃类型（$p = 0.2472$, V = 0.1414）及颜色（$p = 0.6337$）的关联性不显著；基于多元线性回归的风化前成分预测模型在训练集上平均$R^2$达到1.0。问题二构建了决策树分类模型，5折交叉验证平均准确率为86.92\%，标准差为18.78\%；采用K-means聚类进行亚类划分，高钾玻璃与铅钡玻璃的最佳亚类数分别为6类（轮廓系数0.2853）和7类（轮廓系数0.3225）。问题三利用该分类模型对8件未知文物进行鉴别，5件被判定为铅钡玻璃，3件为高钾玻璃。问题四基于皮尔逊相关系数构建成分关联网络，并采用Fisher Z检验量化差异；铅钡玻璃内部存在11条强相关边，平均聚类系数为0.4303，其中PbO与SiO$_2$呈极强的负相关关系（$r = -0.9764$）；而高钾玻璃内部无强相关边，平均聚类系数为0.0，其PbO与SiO$_2$的相关系数为0.0；两类玻璃共有0对化学成分的相关性存在显著差异。

所建模型在风化成分预测（$R^2=1.0$）、玻璃分类（准确率86.92\%）及成分关联分析中均取得有效结果。
\end{abstract}

\textbf{关键词：} 决策树分类；多元线性回归；K-means聚类；关联网络；皮尔逊相关系数；Fisher Z检验；卡方检验